
\section{Poisson Martingale Approximation}

We convert \(W^{\text{RR}}\) into a martingale plus an Abel remainder via the
Poisson equation.
% The variance comparison of the previous section identifies the limiting
% covariance of \(W^{\text{RR}}\), but it does not by itself produce a martingale. The
% deterministic kernel bounds of Sections 4.3--4.4 enter once and for all in
% the boundary/Abel control of the remainder.

\textbf{Poisson kernel.} Let \(\mathsf{Q}\) be the one-step transition kernel:
\begin{equation}
\begin{aligned}
(\mathsf{Q} f)(z) = \mathbb{E} \left[f(Z_{k + 1}) | Z_k = z\right].
\end{aligned}
\end{equation}
Under UGE 1,
\begin{equation}
\begin{aligned}
\lVert  \mathsf{Q}^k \varepsilon  \rVert_\infty \le 2 \lVert  \varepsilon  \rVert_\infty (1 / 4)^{\left\lfloor k / t_{\text{mix}} \right\rfloor}
\end{aligned}
\end{equation}
(valid for centered \(\varepsilon\), \(\pi(\varepsilon) = 0\)) makes the Poisson
series
\begin{equation}
\begin{aligned}
\widehat{\varepsilon} \coloneqq \sum_{k = 0}^\infty \mathsf{Q}^k \varepsilon
\end{aligned}
\end{equation}
absolutely convergent in sup-norm:
\begin{equation}
\begin{aligned}
\lVert  \widehat{\varepsilon}  \rVert_\infty
  \le 2 \lVert  \varepsilon  \rVert_\infty \sum_{k = 0}^\infty (1 / 4)^{\left\lfloor k / t_{\text{mix}} \right\rfloor}
  \le 3 \, t_{\text{mix}} \, \lVert  \varepsilon  \rVert_\infty,
\quad
\lVert  \mathsf{Q} \widehat{\varepsilon}  \rVert_\infty
  \le \lVert  \widehat{\varepsilon}  \rVert_\infty,
\end{aligned}
\end{equation}
and \(\widehat{\varepsilon}\) solves
% The second inequality uses that \(\mathsf{Q}\) is a Markov kernel and contracts
% the sup-norm.
\begin{equation}
\begin{aligned}
\widehat{\varepsilon} - \mathsf{Q} \widehat{\varepsilon} = \varepsilon.
\end{aligned}
\label{eq:poisson-eq}
\end{equation}

\textbf{Conditional centering.} Let \(\mathcal{F}_l \coloneqq \sigma(Z_0, \dots, Z_l)\). For \(l \ge 2\),
\begin{equation}
\begin{aligned}
\varepsilon(Z_l)
  = m_l + t_l,
\quad l \ge 2,
\end{aligned}
\end{equation}
where
\begin{equation}
\begin{aligned}
m_l \coloneqq \widehat{\varepsilon}(Z_l) - \mathsf{Q} \widehat{\varepsilon}(Z_{l - 1}),
\quad
t_l \coloneqq \mathsf{Q} \widehat{\varepsilon}(Z_{l - 1}) - \mathsf{Q} \widehat{\varepsilon}(Z_l),
\quad l \ge 2,
\end{aligned}
\end{equation}
and \(m_l\) is centered conditionally on \(\mathcal{F}_{l - 1}\).
% The Markov property gives
% \(\mathbb{E}[\widehat{\varepsilon}(Z_l)\mid\mathcal{F}_{l-1}]
% =\mathsf{Q}\widehat{\varepsilon}(Z_{l-1})\).
% The \(l = 1\) term is kept as the left Abel boundary rather than included in
% \(M_n^{\text{RR}}\). Abel summation then leaves the discrete derivative
% \(\mathcal{Q}_{l + 1}^{\text{RR}} - \mathcal{Q}_l^{\text{RR}}\), bounded in Section 4.4.

\begin{lemma}\label{lem:poisson-martingale-decomp}
  \textbf{(Stationary Poisson martingale decomposition.)}
  Assume \textbf{UGE 1} and \(\pi(\varepsilon) = 0\). Set
\begin{equation}
\begin{aligned}
  \Delta M_l^{\text{RR}}
    \coloneqq \mathcal{Q}_l^{\text{RR}} \, (\widehat{\varepsilon}(Z_l) - \mathsf{Q} \widehat{\varepsilon}(Z_{l - 1})),
  \quad 2 \le l \le n - 1,
\end{aligned}
\end{equation}
  and let \(M_n^{\text{RR}} \coloneqq \sum_{l = 2}^{n - 1} \Delta M_l^{\text{RR}}\). Then
  \(\left\{\Delta M_l^{\text{RR}}\right\}_{l = 2}^{n - 1}\) is a sequence of \(\mathcal{F}_l\)-martingale
  differences, and
\begin{equation}
\begin{aligned}
  W^{\text{RR}}
    = -\frac{1}{\sqrt{n}} M_n^{\text{RR}} + D_{2, n}^{\text{RR}},
\end{aligned}
\label{eq:poisson-decomp}
\end{equation}
  with the \textbf{Abel remainder}
\begin{equation}
\begin{aligned}
  D_{2, n}^{\text{RR}}
    \coloneqq -\frac{1}{\sqrt{n}} \left[
        \mathcal{Q}_1^{\text{RR}} \, \widehat{\varepsilon}(Z_1)
        + \sum_{l = 1}^{n - 2}
            (\mathcal{Q}_{l + 1}^{\text{RR}} - \mathcal{Q}_l^{\text{RR}}) \, \mathsf{Q} \widehat{\varepsilon}(Z_l)
      \right].
\end{aligned}
\end{equation}
  Moreover, with
% The right boundary
% \(\mathcal{Q}_{n - 1}^{\text{RR}} \, \mathsf{Q} \widehat{\varepsilon}(Z_{n - 1})\)
% vanishes identically because \(\mathcal{Q}_{n - 1}^{\text{RR}} = 0\).
  \(C_{\mathcal{Q}} \coloneqq \lVert  \overline{A}^{-1}  \rVert + 3 C_Q\) a uniform bound on
  \(\lVert  \mathcal{Q}_l^{\text{RR}}  \rVert\), and \(C_2\) the constant from the Corollary of
  Section 4.4,
\begin{equation}
\begin{aligned}
  \lVert  D_{2, n}^{\text{RR}}  \rVert_\infty
    \le \frac{3 \, t_{\text{mix}} \, \lVert  \varepsilon  \rVert_\infty}{\sqrt{n}}
       \left(C_{\mathcal{Q}} + \frac{C_2}{a^2}\right).
\end{aligned}
\label{eq:D2-bound}
\end{equation}
  Consequently, for every \(p \ge 1\),
\begin{equation}
\begin{aligned}
  \lVert  D_{2, n}^{\text{RR}}  \rVert_{L_p}
    \le \frac{C \, t_{\text{mix}} \, \lVert  \varepsilon  \rVert_\infty}{a^2 \, \sqrt{n}},
\end{aligned}
\end{equation}
  with a constant \(C\) depending only on \(\lVert  \overline{A}^{-1}  \rVert\), \(\kappa_Q\),
  and \(\lVert  \overline{A}  \rVert\).
\end{lemma}

\textit{Proof.} The increments are martingale differences by the Markov property.
% Here \(\mathcal{Q}_l^{\text{RR}}\) is deterministic and
% \(\mathbb{E}[\widehat{\varepsilon}(Z_l)\mid\mathcal{F}_{l-1}]
% =\mathsf{Q}\widehat{\varepsilon}(Z_{l-1})\).

Substitute \eqref{eq:poisson-eq} into \eqref{eq:W-RR}:
\begin{equation}
\begin{aligned}
W^{\text{RR}}
  = -\frac{1}{\sqrt{n}} \sum_{l = 1}^{n - 1}
    \mathcal{Q}_l^{\text{RR}} \, (\widehat{\varepsilon}(Z_l) - \mathsf{Q} \widehat{\varepsilon}(Z_l)).
\end{aligned}
\end{equation}
Peel off \(l = 1\):
\begin{equation}
\begin{aligned}
W^{\text{RR}}
  = -\frac{1}{\sqrt{n}} \mathcal{Q}_1^{\text{RR}} \, \widehat{\varepsilon}(Z_1)
    + \frac{1}{\sqrt{n}} \mathcal{Q}_1^{\text{RR}} \, \mathsf{Q} \widehat{\varepsilon}(Z_1)
    - \frac{1}{\sqrt{n}} M_n^{\text{RR}}
    - \frac{1}{\sqrt{n}} \sum_{l = 2}^{n - 1}
      \mathcal{Q}_l^{\text{RR}} \, (\mathsf{Q} \widehat{\varepsilon}(Z_{l - 1}) - \mathsf{Q} \widehat{\varepsilon}(Z_l)).
\end{aligned}
\end{equation}
Set \(g_l \coloneqq \mathsf{Q} \widehat{\varepsilon}(Z_l)\). Abel summation gives
\begin{equation}
\begin{aligned}
\sum_{l = 2}^{n - 1} \mathcal{Q}_l^{\text{RR}} (g_{l - 1} - g_l)
  = \mathcal{Q}_2^{\text{RR}} \, g_1
    - \mathcal{Q}_{n - 1}^{\text{RR}} \, g_{n - 1}
    + \sum_{l = 2}^{n - 2}
      (\mathcal{Q}_{l + 1}^{\text{RR}} - \mathcal{Q}_l^{\text{RR}}) \, g_l.
\end{aligned}
\end{equation}
The boundary terms combine as
\begin{equation}
\begin{aligned}
&-\mathcal{Q}_1^{\text{RR}} \, \widehat{\varepsilon}(Z_1)
+ \mathcal{Q}_1^{\text{RR}} \, g_1
- \mathcal{Q}_2^{\text{RR}} \, g_1
+ \mathcal{Q}_{n - 1}^{\text{RR}} \, g_{n - 1}
- \sum_{l = 2}^{n - 2}
    (\mathcal{Q}_{l + 1}^{\text{RR}} - \mathcal{Q}_l^{\text{RR}}) \, g_l \\
&\quad= -\mathcal{Q}_1^{\text{RR}} \, \widehat{\varepsilon}(Z_1)
       - \sum_{l = 1}^{n - 2}
          (\mathcal{Q}_{l + 1}^{\text{RR}} - \mathcal{Q}_l^{\text{RR}}) \, g_l
       + \mathcal{Q}_{n - 1}^{\text{RR}} \, g_{n - 1}.
\end{aligned}
\end{equation}
The rightmost boundary vanishes because
\(\mathcal{Q}_{n - 1}^{\text{RR}} = 2 \alpha I - 2 \alpha I = 0\). This gives the stated
formula.
% This uses \(Q_{n - 1}^{(\alpha)} = \alpha I\) and
% \(Q_{n - 1}^{(2 \alpha)} = 2 \alpha I\); the sign is absorbed into
% \(D_{2, n}^{\text{RR}}\).

For \eqref{eq:D2-bound}, use the sup-norm Poisson bound, the uniform weight bound, and
the summed-total-variation bound:
\begin{equation}
\begin{aligned}
\lVert  \sqrt{n} \, D_{2, n}^{\text{RR}}  \rVert
  \le 3 t_{\text{mix}} \lVert  \varepsilon  \rVert_\infty (C_{\mathcal{Q}} + C_2 / a^2).
\end{aligned}
\end{equation}
The \(L_p\) bound is immediate from the deterministic sup-norm bound. \(\square\)
